% =====================================================================
% Kapitel 2: Grundlagen
% =====================================================================

\chapter{Grundlagen}

\label{ch:grundlagen}

Dieses Kapitel erklärt alle Fachbegriffe, die in den Folgekapiteln
verwendet werden. Wer mit allen Konzepten vertraut ist, kann es
überspringen und bei Unklarheiten als Nachschlagewerk nutzen.

\section{Datenbank-Grundlagen}

\begin{definition}
\textbf{Relationale Datenbank}
Eine \emph{relationale Datenbank} organisiert Daten in
\emph{Tabellen} (auch \emph{Relationen} genannt) mit Zeilen (Datensätze)
und Spalten (Attribute). Zwischen den Tabellen bestehen
\emph{Beziehungen} über \emph{Fremdschlüssel} (\foreign{foreign key}).
Die meisten modernen Datenbanken -- einschließlich PostgreSQL --
sprechen die standardisierte Abfragesprache SQL.
\end{definition}

\begin{definition}
\textbf{SQL (Structured Query Language)}
SQL ist die standardisierte Abfragesprache für relationale Datenbanken.
SQL unterteilt sich in Untergruppen:

\begin{itemize}
  \item \textbf{DDL} (Data Definition Language): \code{CREATE TABLE},
        \code{ALTER TABLE}, \code{DROP TABLE}.
  \item \textbf{DML} (Data Manipulation Language): \code{SELECT},
        \code{INSERT}, \code{UPDATE}, \code{DELETE}.
  \item \textbf{DCL} (Data Control Language): \code{GRANT},
        \code{REVOKE}.
\end{itemize}
\end{definition}

\begin{definition}
\textbf{PostgreSQL}
PostgreSQL -- oft kurz \emph{Postgres} -- ist ein
\emph{objekt-relationales} Datenbankmanagementsystem (ORDBMS). Es ist
quelloffen, kostenlos, seit über 30~Jahren in Entwicklung und gilt als
eine der stabilsten und funktionsreichsten Open-Source-Datenbanken.
PostgreSQL unterstützt JSON/JSONB, Volltextsuche, Window Functions und
vieles mehr.
\end{definition}

\begin{beispiel}
\textbf{PostgreSQL-Verbindung in Python}
\begin{lstlisting}[language=python]
import psycopg

conn = psycopg.connect(
    host="localhost",
    port=5432,
    user="deephold",
    password="deephold",
    dbname="deephold",
)
with conn.cursor() as cur:
    cur.execute("SELECT COUNT(*) FROM macro_observations")
    print(cur.fetchone())
\end{lstlisting}
\end{beispiel}

\begin{definition}
\textbf{Index}
Ein \emph{Index} ist eine zusätzliche Datenstruktur, die das schnelle
Suchen in einer Tabelle ermöglicht -- analog zum Stichwortverzeichnis
eines Buchs. Ein Index auf Spalte~\code{X} beschleunigt Abfragen
mit \code{WHERE X = ...}, verlangsamt aber \code{INSERT}-Operationen
und kostet zusätzlichen Speicherplatz. Daher werden nur Spalten
indiziert, die häufig in \code{WHERE}-Klauseln vorkommen.
\end{definition}

\begin{beispiel}
\textbf{Index-Definition}
\begin{lstlisting}[language=sql]
CREATE INDEX ix_prices_daily_date
  ON prices_daily (date);

CREATE UNIQUE INDEX ix_instrument_identifiers_scheme_value
  ON instrument_identifiers (scheme, value);
\end{lstlisting}
\end{beispiel}

\begin{definition}
\textbf{JSONB}
JSONB ist die binäre Variante des JSON-Datentyps in PostgreSQL. Werte
werden nicht als Text, sondern als effizient geparste interne Struktur
gespeichert. Das ermöglicht Indexierung auf JSON-Feldern und
schnelle Pfad-Queries (\code{data-\texttt{>}'field'}).
\end{definition}

\section{Python-Grundlagen}

\begin{definition}
\textbf{Python}
Python ist eine interpretierte, dynamisch typisierte
Hochsprachen-Programmiersprache. Sie wird im Data-Science-Umfeld
besonders wegen ihrer einfachen Syntax, ihrer umfangreichen
Standardbibliothek und ihres dichten Ökosystems an
Dritt-Bibliotheken geschätzt. Das Projekt verwendet Python in
Version~3.11 oder neuer.
\end{definition}

\begin{definition}
\textbf{Poetry}
Poetry ist ein modernes Tool für \emph{Dependency-Management} in
Python-Projekten. Es verwaltet die \code{pyproject.toml}-Datei, legt
virtuelle Environments automatisch an und löst transitive
Abhängigkeiten auf. In diesem Projekt wird es durch
\code{pip} und ein venv-Verzeichnis ersetzt (siehe Kapitel~\ref{ch:infrastruktur}),
weil Poetry nicht überall verfügbar ist.
\end{definition}

\begin{definition}
\textbf{Pandas}
Pandas ist die Standard-Bibliothek für tabellarische Datenverarbeitung
in Python. Die zentrale Datenstruktur ist das \emph{DataFrame} -- eine
beschriftete, zweidimensionale Tabelle. Pandas wird für Daten-
Cleaning, Resampling, Grouping und vieles mehr verwendet.
\end{definition}

\begin{definition}
\textbf{Polars}
Polars ist eine neuere, schnellere Alternative zu Pandas, implementiert
in Rust. Polars nutzt \emph{Lazy Evaluation} und \emph{Columnar
Storage}, was es auf großen Datensätzen drastisch schneller macht.
Die API ähnelt Pandas, ist aber stärker typisiert.
\end{definition}

\begin{beispiel}
\textbf{Polars vs.\ Pandas -- ein einfaches Beispiel}
\begin{lstlisting}[language=python]
import polars as pl
import pandas as pd

# Beide erzeugen das gleiche Ergebnis:
df_pl = pl.DataFrame({"a": [1, 2, 3], "b": [4.0, 5.0, 6.0]})
df_pd = pd.DataFrame({"a": [1, 2, 3], "b": [4.0, 5.0, 6.0]})

# Polars:
mean_b = df_pl["b"].mean()  # Lazy, gibt eine Expression zurück
print(df_pl.select(pl.col("b").mean()))  # 5.0

# Pandas:
mean_b = df_pd["b"].mean()  # Eager, gibt einen float zurück
print(mean_b)  # 5.0
\end{lstlisting}
\end{beispiel}

\begin{definition}
\textbf{SQLAlchemy}
SQLAlchemy ist ein \emph{ORM-Toolkit} (Object-Relational Mapper) für
Python. Es erlaubt, Datenbank-Tabellen als Python-Klassen zu
modellieren und Abfragen in Python-Syntax zu schreiben, ohne rohes
SQL zu schreiben. SQLAlchemy unterstützt zwei Stile:

\begin{itemize}
  \item \textbf{Core}: niedrigschwellige SQL-Bausteine
        (\code{select()}, \code{insert()} etc.).
  \item \textbf{ORM}: hochschwellige, objekt-orientierte
        Zugriffsschicht mit Klassen für jede Tabelle.
\end{itemize}
In diesem Projekt wird die ORM-Schicht verwendet.
\end{definition}

\begin{beispiel}
\textbf{SQLAlchemy-Definition einer Tabelle}
\begin{lstlisting}[language=python]
from sqlalchemy.orm import DeclarativeBase, Mapped, mapped_column

class Base(DeclarativeBase):
    pass

class Instrument(Base):
    __tablename__ = "instruments"
    instrument_id: Mapped[int] = mapped_column(primary_key=True)
    asset_class: Mapped[str] = mapped_column()
    name: Mapped[str] = mapped_column()
    # ...
\end{lstlisting}
\end{beispiel}

\begin{definition}
\textbf{Alembic}
Alembic ist das offizielle Migrationswerkzeug für SQLAlchemy. Es
verwaltet \emph{Schema-Migrationen} -- das Hinzufügen, Ändern oder
Löschen von Tabellen und Spalten, ohne Daten zu verlieren. Jede
Migration ist eine Python-Datei mit zwei Funktionen: \code{upgrade()}
und \code{downgrade()}.
\end{definition}

\begin{beispiel}
\textbf{Alembic-Migration -- sehr vereinfacht}
\begin{lstlisting}[language=python]
def upgrade() -> None:
    op.create_table(
        "macro_series",
        sa.Column("series_id", sa.Text, primary_key=True),
        sa.Column("name", sa.Text),
        # ...
    )

def downgrade() -> None:
    op.drop_table("macro_series")
\end{lstlisting}
\end{beispiel}

\section{Docker und Container}

\begin{definition}
\textbf{Container}
Ein \emph{Container} ist eine isolierte Laufzeit-Umgebung für eine
Anwendung. Anders als bei virtuellen Maschinen wird nicht der
gesamte Rechner virtualisiert, sondern nur die
\emph{Anwendungs-Layer} (Binaries, Libraries, Konfiguration).
Container starten in Sekunden, nicht in Minuten, und sind sehr
ressourcenschonend.
\end{definition}

\begin{definition}
\textbf{Docker}
Docker ist die populärste Container-Engine. Sie packt eine Anwendung
in ein \emph{Image} (eine Art Bauplan) und startet daraus
\emph{Container}. Ein Image wird durch ein \code{Dockerfile}
beschrieben, ein Container ist eine laufende Instanz eines Images.
\end{definition}

\begin{definition}
\textbf{Docker Compose}
Docker Compose ist ein Werkzeug zur Definition und zum Start
\emph{mehrerer} Container als zusammengehörige Anwendung. Die
Konfiguration erfolgt in einer \code{docker-compose.yml}-Datei.
Mit \code{docker compose up -d} startet man alle Services im
Hintergrund; mit \code{docker compose down} stoppt man sie wieder.
\end{definition}

\begin{beispiel}
\textbf{docker-compose.yml -- Auszug}
\begin{lstlisting}[language=yaml]
services:
  postgres:
    image: postgres:16-alpine
    environment:
      POSTGRES_USER: deephold
      POSTGRES_PASSWORD: deephold
      POSTGRES_DB: deephold
    ports:
      - "5432:5432"
    healthcheck:
      test: ["CMD-SHELL", "pg_isready -U deephold"]
      interval: 5s
\end{lstlisting}
\end{beispiel}

\begin{definition}
\textbf{Docker Volume}
Ein \emph{Volume} ist ein persistenter Speicherbereich, der vom
Container unabhängig existiert. Wird ein Container gelöscht, bleiben
Daten im Volume erhalten. \code{deephold\_db} definiert ein Volume
\code{postgres\_data} für die Datenbank, damit die Daten
Container-Updates überleben.
\end{definition}

\section{Networking-Grundlagen}

\begin{definition}
\textbf{HTTP (Hypertext Transfer Protocol)}
HTTP ist das Standard-Protokoll für die Kommunikation zwischen
Web-Clients und Web-Servern. Datenquellen wie FRED und Yahoo Finance
bieten ihre Daten über \emph{RESTful HTTP}-APIs an. Ein typischer
Aufruf ist \code{GET
https://api.example.com/data?param=value}, der mit einem
JSON-Array oder XML-Dokument beantwortet wird.
\end{definition}

\begin{definition}
\textbf{Rate Limiting}
Die meisten öffentlichen APIs begrenzen, wie viele Anfragen ein
einzelner Client pro Zeiteinheit stellen darf -- das sogenannte
\emph{Rate Limiting}. Wird das Limit überschritten, antwortet der
Server mit HTTP 429 (Too Many Requests). Das Projekt implementiert
ein \emph{Token-Bucket}-Verfahren, bei dem jeder Vendor ein eigenes
Bucket mit konfigurierbarer Rate und Burst-Größe hat.
\end{definition}

\begin{beispiel}
\textbf{Token-Bucket in Python}
\begin{lstlisting}[language=python]
import time

class TokenBucket:
    def __init__(self, rate_per_min: int, burst: int):
        self.rate = rate_per_min / 60.0   # tokens per second
        self.burst = burst
        self.tokens = float(burst)
        self.last = time.monotonic()

    def take(self, n: int = 1) -> None:
        now = time.monotonic()
        elapsed = now - self.last
        self.last = now
        self.tokens = min(self.burst, self.tokens + elapsed * self.rate)
        if self.tokens < n:
            # Blockierend warten
            time.sleep((n - self.tokens) / self.rate)
            self.tokens = 0
        else:
            self.tokens -= n
\end{lstlisting}
\end{beispiel}

\section{TypeScript- und TUI-Grundlagen}

\begin{definition}
\textbf{TypeScript}
TypeScript ist eine Obermenge von JavaScript, die statische Typen
hinzufügt. TypeScript-Code wird zu JavaScript compiliert und läuft
dann in jeder JavaScript-Engine (z.\,B.\ Node.js oder Bun).
\end{definition}

\begin{definition}
\textbf{Bun}
Bun ist eine moderne JavaScript-/TypeScript-Runtime, ähnlich wie
Node.js, aber deutlich schneller beim Start und bei der Installation
von Dependencies. Bun enthält einen eingebauten
\emph{Paketmanager}, einen \emph{Test-Runner} und einen
\emph{Bundler}. Bun ist die Standard-Runtime für das in diesem
Projekt verwendete OpenTUI-Framework.
\end{definition}

\begin{definition}
\textbf{Terminal User Interface (TUI)}
Eine \emph{TUI} ist ein textbasiertes Interface, das im Terminal
läuft -- im Gegensatz zu einem GUI (grafisches User Interface).
Im Gegensatz zu einem Kommandozeilen-Tool, das nur Text ein- und
ausgibt, bietet eine TUI ein \emph{interaktives} Layout mit
\emph{Widgets} wie Listen, Buttons, Formularfeldern und Charts.
Beispiele: vim, htop, lazygit.
\end{definition}

\begin{definition}
\textbf{OpenTUI}
OpenTUI ist ein von Anomaly (https://anoma.ly) entwickeltes
TUI-Framework. Es besteht aus einem in Zig geschriebenen
\emph{Core} und TypeScript-Bindings mit React- und Solid.js-Support.
Der Kern nutzt Yoga (Facebooks Flexbox-Engine), um Layouts zu
berechnen -- genau wie im Web -- und rendert das Ergebnis in einen
Terminal-Buffer.
\end{definition}

\begin{definition}
\textbf{Yoga (Flexbox-Layout)}
Yoga ist Facebooks plattformübergreifender Layout-Algorithmus, der
\foreign{flexbox} -- das CSS-Layout-System -- in native Bibliotheken
portiert. OpenTUI nutzt Yoga, um komplexe Layouts (Sidebars,
Header, Footer, mehrspaltige Inhalte) pixelgenau im Terminal zu
positionieren.
\end{definition}

\section{Finanzmarkt-Grundlagen}

\begin{definition}
\textbf{Adjusted Close}
Der \emph{adjusted Close} ist der Schlusskurs eines Wertpapiers,
\emph{nachbereinigt} um alle Kapitalmaßnahmen, die zwischen
Kaufzeitpunkt und aktuellem Zeitpunkt stattgefunden haben -- vor
allem Splits und Dividenden. Wenn eine Aktie am Tag~X zu 100\,{}
schließt und am Folgetag einen 2:1-Split macht, dann steht der
adjusted Close für Tag~X nicht bei 100, sondern bei 50\,{}.
Backtests, die adjusted close verwenden, sind -- in der Retail-Praxis
-- korrekt für historische Investoren.
\end{definition}

\begin{definition}
\textbf{Log-Return vs.\ Simple Return}
Es gibt zwei gängige Arten, Renditen zu berechnen:

\begin{itemize}
  \item \textbf{Simple Return}: $R_t = \frac{P_t - P_{t-1}}{P_{t-1}}$
  \item \textbf{Log-Return}: $r_t = \ln\frac{P_t}{P_{t-1}}$
\end{itemize}

Log-Returns sind \emph{zeitadditiv} -- der Log-Return über zwei
Tage ist die Summe der täglichen Log-Returns. Dies vereinfacht
viele Berechnungen. Bei kleinen Returns (~$|r_t| < 5\%$) sind die
beide Werte praktisch identisch.
\end{definition}

\begin{definition}
\textbf{Volatilität (Vol)}
Die \emph{Volatilität} ist die annualisierte Standardabweichung der
Renditen. Sie ist das Standardmaß für das Risiko eines Wertpapiers
oder Portfolios. Übliche Symbole: $\sigma$ (Sigma) oder $\text{vol}$.
\end{definition}

\begin{beispiel}
\textbf{Volatilität berechnen}
\begin{lstlisting}[language=python]
import math
import statistics

# Tägliche Log-Returns einer Aktie
daily_returns = [0.01, -0.02, 0.005, 0.015, -0.01, 0.0, 0.02]

# Standardabweichung der täglichen Returns
daily_vol = statistics.stdev(daily_returns)  # ≈ 0.0124

# Annualisiert (252 Handelstage)
annual_vol = daily_vol * math.sqrt(252)     # ≈ 0.196 = 19.6%
\end{lstlisting}
\end{beispiel}

\begin{definition}
\textbf{CAGR (Compound Annual Growth Rate)}
Die \emph{CAGR} ist die durchschnittliche jährliche Wachstumsrate
eines Investments über einen bestimmten Zeitraum, unter Annahme
eines geometrischen (zinseszinsartigen) Wachstums:

\[
  \text{CAGR} = \left(\frac{V_\text{end}}{V_\text{start}}\right)^{1/n} - 1
\]

wobei $n$ die Anzahl der Jahre ist. Ein Investment, das von 100\,{}
auf 200\,{} in 5~Jahren wächst, hat eine CAGR von
$\sqrt[5]{200/100} - 1 \approx 14{,}87\%$.
\end{definition}

\begin{definition}
\textbf{Sharpe Ratio}
Die \emph{Sharpe Ratio} misst die \emph{renditebereinigte Performance}
eines Investments -- die Überrendite pro Einheit Gesamtrisiko:

\[
  \text{Sharpe} = \frac{\bar{R}_p - R_f}{\sigma_p}
\]

wobei $\bar{R}_p$ die annualisierte Portfoliorendite, $R_f$ der
risikofreie Zins (oft 4\,\% p.\,a.) und $\sigma_p$ die annualisierte
Volatilität ist. Höher ist besser; eine Sharpe~>~1 gilt als gut,
Sharpe~>~2 als exzellent.
\end{definition}

\begin{definition}
\textbf{Sortino Ratio}
Eine Variante der Sharpe Ratio, die nur die
\emph{Abwärtsvolatilität} (negative Renditen) als Risikomaß
verwendet. Sie ist aussagekräftiger für asymmetrische Strategien, die
gerne hohe positive Volatilität (große Gewinne) in Kauf nehmen, aber
negative Volatilität (Verluste) minimieren wollen.
\end{definition}

\begin{definition}
\textbf{Maximum Drawdown (MaxDD)}
Der \emph{Maximum Drawdown} ist der größte prozentuale Verlust vom
einem vorherigen Höchststand bis zu einem nachfolgenden Tiefststand:

\[
  \text{MaxDD} = \min_t \frac{V_t - V_\text{peak}}{V_\text{peak}}
\]

Ein MaxDD von $-30\,\%$ bedeutet, dass das Portfolio irgendwann
30\,\% unter seinem vorherigen Höchststand lag. Dies ist das
wichtigste \emph{Tail-Risk}-Maß.
\end{definition}

\begin{definition}
\textbf{Walk-Forward-Backtest}
Ein \emph{Walk-Forward-Backtest} ist die realistischste Form der
Backtest-Validierung. Anders als bei einem klassischen
\emph{In-Sample}-Backtest wird das Modell nicht auf denselben Daten
trainiert und evaluiert, auf denen es entwickelt wurde. Stattdessen:

\begin{enumerate}
  \item Trainiere auf $[t - L, t]$ (z.\,B.\ 12~Monate)
  \item Evaluiere auf $[t, t+1]$ (z.\,B.\ 1~Monat)
  \item Wandere ein Fenster weiter, wiederhole.
\end{enumerate}

So bekommt das Modell zu jedem Zeitpunkt nur Daten, die es auch
\emph{tatsächlich} gehabt hätte -- kein \emph{Look-Ahead-Bias}.
\end{definition}

\begin{definition}
\textbf{Look-Ahead-Bias}
\emph{Look-Ahead-Bias} entsteht, wenn ein Backtest Informationen
verwendet, die zum Entscheidungszeitpunkt in der Realität noch nicht
verfügbar gewesen wären -- zum Beispiel ein Split, der erst nach dem
Backtest-Datum in den Preisdaten nachträglich eingerechnet wurde.
Walk-Forward-Tests eliminieren Look-Ahead-Bias vollständig.
\end{definition}

\begin{definition}
\textbf{Survivorship-Bias}
\emph{Survivorship-Bias} ist die Tendenz, nur die
\emph{überlebenden} Unternehmen in einem Index zu betrachten. Wenn
Sie z.\,B.\ die durchschnittliche Rendite der heutigen S\&P-500-Mitglieder
in der Vergangenheit berechnen, ignorieren Sie all die Unternehmen,
die zwischenzeitlich pleitegingen oder delistet wurden. Das
überschätzt die wahre historische Performance.
\end{definition}

\begin{definition}
\textbf{Momentum}
\emph{Momentum} ist die Tendenz von Aktien, die in der Vergangenheit
gut gelaufen sind, dies auch in der nahen Zukunft zu tun -- und
umgekehrt. Die klassische Momentum-Strategie (Jegadeesh \&
Titman, 1993) kauft die Top-10\,\% der zurückliegenden 12~Monate
und leert die unteren 10\,\%. Momentum ist eines der am besten
dokumentierten Anomalien in der Finanzmarktforschung.
\end{definition}

\begin{definition}
\textbf{VAM (Volatility-Adjusted Momentum)}
Die \emph{VAM} (manchmal auch \emph{Sharpe-Momentum} genannt)
normalisiert die historische Rendite durch die historische
Volatilität:

\[
  \text{VAM}_i = \frac{R_i}{\sigma_i}
\]

Eine Aktie mit 50\,\% Jahresrendite und 30\,\% Volatilität hat
einen VAM von $1{,}67$. Eine andere mit 50\,\% Rendite aber
60\,\% Volatilität hat einen VAM von $0{,}83$ -- sie ist bei
gleicher Rendite deutlich riskanter. VAM hilft, den
\emph{Winner's Curse} zu vermeiden: Aktien, die nur wegen hoher
Beta gut laufen, werden aussortiert.
\end{definition}

\begin{definition}
\textbf{Skip-Month}
Beim Walk-Forward-Backtest ist es üblich, den letzten Monat vor
dem Re-Balancing-Datum \emph{auszulassen} (das sogenannte
\emph{Skip-Month}). Grund: Die jüngsten 21~Tage enthalten oft
Mean-Reversion-Effekte und Mikrostruktur-Rauschen, die die
Signale verzerren würden. Das AEGIS-Paper und dieses Projekt
verwenden ein 21-Tage-Skip (Default).
\end{definition}

\begin{definition}
\textbf{Friction / Transaction Costs}
\emph{Friction} (auch \emph{Transaction Costs}) sind die
\emph{Transaktionskosten} eines Backtests -- typischerweise als
Bps (Basispunkte, 1~Bp = 0{,}01\,\%) pro umgesetztem Dollar
ausgedrückt. Das AEGIS-Paper verwendet 10~Bp pro Handel, was
etwa 1~Bp pro Monat-Portfolio-Turnover entspricht.
\end{definition}

\section{Was kommt als Nächstes?}

Mit diesen Grundlagen bewaffnet, können Sie sich auf die
Infrastruktur (Kapitel~\ref{ch:infrastruktur}) stürzen. Dort
sehen wir, wie alle diese Komponenten in einem lauffähigen
System zusammenspielen.
